\documentclass[aspectratio=169]{beamer}
%\documentclass[aspectratio=169,handout]{beamer}

% Trying to use notes on separate page
% \documentclass[aspectratio=169,notes]{beamer}       % print frame + notes
% \documentclass[notes=only]{beamer}   % only notes
% \documentclass{beamer}              % only frames
%\usepackage{pgfpages}
%\setbeameroption{show notes}
%\setbeameroption{show notes on second screen=right}
%

%\setbeamercovered{transparent}

\usetheme[titleformat=regular%
%,numbering=fraction% use slide numbers
]{metropolis}
\metroset{%
  progressbar=foot,%
  %background=dark,
  block=fill
}
\only<handout>{\metroset{sectionpage=none}}
\only<handout>{\usecolortheme{dove}}
\usepackage{appendixnumberbeamer} % separate appendix
\usepackage[citestyle=authortitle,sorting=none]{biblatex}
\setbeamerfont{footnote}{size=\tiny}
\addbibresource{mae.bib}

% multimedia in beamer presentations
\usepackage{multimedia}

% adding some better facilities to latex such as proper booleans
\usepackage{etoolbox}

% Unicode math symbols for XeLaTeX
\usepackage{mathrsfs}

% standard math packages
\usepackage{amssymb}
\usepackage{amsthm}
\usepackage{amsmath}
\usepackage{amstext}
\usepackage{mathabx}
\usepackage{stmaryrd}
% math tools for amsmath
\usepackage{mathtools}

\usepackage{proof}

\usepackage{pgfpages}

\usepackage{listings}
\usepackage{xcolor}

% subfigures, subfloats
\usepackage{subcaption}

% for multiinclude
\usepackage{xmpmulti}

% tikz & friends

\usepackage{galois}
\usepackage{tikz}
\usetikzlibrary{fit,calc,shapes,arrows.meta,patterns,backgrounds}
\usetikzlibrary{decorations.pathmorphing}
\usetikzlibrary{cd}
\usetikzlibrary{positioning}
\usepackage[beamer]{hf-tikz}

% for author-title-year citing
\usepackage{xpatch}

\newcommand{\tikzmark}[1]{%
  \tikz[overlay,remember picture,baseline] \node [anchor=base] (#1) {};}

\newenvironment{tightcenter}{%
  \setlength\topsep{0pt}
  \setlength\parskip{0pt}
  \begin{center}
}{%
  \end{center}
}

% new colors
\definecolor{lightblue}{RGB}{217,220,253}
\definecolor{lightred}{RGB}{251,216,218}

\DeclareMathOperator{\pw}{\mathcal{P}} % powerset
\newcommand{\fset}[1]{\mathsf{#1}}
\newcommand{\nats}{\mathbb{N}}
\newcommand{\zahlen}{\mathbb{Z}}
\newcommand{\bools}{\mathbb{B}}
\newcommand{\Set}[1]{\left\{#1\right\}}
\newcommand{\true}{\kw{true}}
\newcommand{\false}{\kw{false}}
\newcommand{\sidenote}[1]{\hfill\quad \textsf{#1}}

\newcommand{\disunion}{+}

% named set
\newcommand{\ns}[1]{\mathit{#1}}
% function
\newcommand{\fn}[1]{\mathrm{#1}}
% "vector of"
\newcommand{\vo}[1]{\overrightarrow{#1}}
% "set of"
\newcommand{\setOf}[1]{\overline{#1}}
% syntactic tag
\newcommand{\sTag}[2]{\textsf{\textbf{#1}}\,#2}
% keyword
\newcommand{\kw}[1]{\texttt{#1}}
% usual suspects
\newcommand{\State}{\ns{State}}
\newcommand{\Value}{\ns{Value}}
\newcommand{\Stmt}{\ns{Stmt}}
\newcommand{\Env}{\ns{Env}}
\newcommand{\Store}{\ns{Store}}
\newcommand{\Kont}{\ns{Kont}}
% syntactic domains
\newcommand{\Exp}{\ns{Exp}}
\newcommand{\Var}{\ns{Var}}
\newcommand{\Addr}{\ns{Addr}}
% put a value to a pointer
\newcommand{\update}{\leftarrow}

% ebnf
\newcommand{\eDEF}{\,::=\;}
\newcommand{\eOR}{\;\vert\;}

\newcommand{\widen}{\nabla}

% \newcommand{\|}{\,\vert\,}

\newcommand{\todo}[1]{\iftoggle{TODO}{\textcolor{red}{TODO: #1}}{}}
% ceiling and floor symbols
\DeclarePairedDelimiter\ceil{\lceil}{\rceil}
\DeclarePairedDelimiter\floor{\lfloor}{\rfloor}

% big O notation
\DeclareMathOperator{\bigO}{O}

% fixed points
\DeclareMathOperator{\lfp}{lfp}

% print both years for bibliography
\renewbibmacro*{cite:labelyear+extrayear}{%
\iffieldundef{labelyear}
{}
{\printtext[bibhyperref]{%
\iffieldundef{origyear}{}{\printfield{origyear}\addslash}%   <--- added
\printfield{labelyear}%
\printfield{extrayear}}}}

\renewbibmacro*{date+extrayear}{%
\iffieldundef{year}
{}
{\printtext[parens]{%
\iffieldundef{origyear}{}{\printfield{origyear}\addslash}%  <--- added
\printdateextra}}}

% overlay an image
\def\Put(#1,#2)#3{\leavevmode\makebox(0,0){\put(#1,#2){#3}}}

% text over symbols nicely, requires amsmath for overset
\newcommand\textoverop[2]{\mathrel{\overset{\makebox[0pt]{\mbox{\normalfont\tiny\sffamily #1}}}{#2}}}

% special arrows
\newcommand\monarrow{\textoverop{mon}{\rightarrow}}

% theorems
\newtheorem{thm}{Theorem}
\newtheorem{eg}{Example}

\newcommand{\abscolor}[1]{\textcolor{mLightBrown}{#1}}
\newcommand{\concolor}[1]{\textcolor{mLightGreen}{#1}}
\newcommand{\abst}[1]{#1^{\#}}

\newcommand{\step}{\rightsquigarrow}

\newcommand{\altm}{\; {\color{black}\mid}\; }

% listings setup
%\lstset{basicstyle=\tiny\ttfamily,columns=fixed}

\usepackage{fnpct}
\AdaptNoteOpt\footcite\multfootcite

\xapptobibmacro{cite}{\setunit{\nametitledelim}\printfield{year}}{}{}

% highlighting lines in listings

% two approaches: https://tex.stackexchange.com/questions/8851/how-can-i-highlight-some-lines-from-source-code)
\lstset{
  basicstyle=\scriptsize\ttfamily,language=[LaTeX]Tex,breaklines=true,
  breakautoindent=true,breakindent=2ex,
}

\tikzset{onslide/.code args={<#1>#2}{%
  \only<#1>{\pgfkeysalso{#2}} % \pgfkeysalso doesn't change the path
}}

\makeatletter
\newenvironment<>{btHighlight}[1][]
{\begin{onlyenv}#2\begingroup\tikzset{bt@Highlight@par/.style={#1}}\begin{lrbox}{\@tempboxa}}
{\end{lrbox}\bt@HL@box[bt@Highlight@par]{\@tempboxa}\endgroup\end{onlyenv}}

\newcommand<>\btHL[1][]{%
  \only#2{\begin{btHighlight}[#1]\bgroup\aftergroup\bt@HL@endenv}%
}
\def\bt@HL@endenv{%
  \end{btHighlight}%   
  \egroup
}
\newcommand{\bt@HL@box}[2][]{%
  \tikz[#1]{%
    \pgfpathrectangle{\pgfpoint{1pt}{0pt}}{\pgfpoint{\wd #2}{\ht #2}}%
    \pgfusepath{use as bounding box}%
    \node[anchor=base west, fill=orange!30,outer sep=0pt,inner xsep=1pt, inner ysep=0pt, rounded corners=3pt, minimum height=\ht\strutbox+1pt,#1]{\raisebox{1pt}{\strut}\strut\usebox{#2}};
  }%
}
\makeatother

% more coloring in equations?
\newcommand{\colorIt}[3]{
   \tikz[baseline]{ \node[fill=#1!20,anchor=base,rounded corners=2pt] (#2) {\ensuremath{#3}};}
}

\usepackage{listings-rust}

%%% Local Variables:
%%% mode: latex
%%% TeX-master: "main"
%%% TeX-engine: xetex
%%% End:


\only<handout>{
  \pgfpagesuselayout{4 on 1}[letterpaper,border shrink=5mm,landscape]
}

\newtoggle{notes}
%\only<beamer>{\toggletrue{notes}}

% add notes:
\iftoggle{notes}{
  \makeatletter
  \def\beamer@framenotesbegin{% at beginning of slide
    %\gdef\beamer@noteitems{}%
    %\gdef\beamer@notes{}%
    \usebeamercolor[fg]{normal text}
  }
  \makeatother
  \setbeamertemplate{note page}[plain]
  \setbeamerfont{note page}{size=\footnotesize}
  \setbeameroption{show notes on second screen=right}
}{}

%\newtoggle{labdemo}
%\toggletrue{labdemo}
\newtoggle{TODO}
\toggletrue{TODO}

\title[Major Area Exam]{Memory Safety in Systems Languages} %Techniques for
%\subtitle{Tradeoffs in Efficiency and Completeness}
%\subtitle{A Balancing Act}
\subtitle{Major Area Exam}
\date{June 11, 2018}
\author{Michael Christensen}
\institute[UCSB]{
  \normalsize
  {\large \bfseries Committee:}\\
  Ben Hardekopf (\,$\vcenter{\hbox{\includegraphics[height=1em]{chair/file.eps}}}$) \quad
  Tim Sherwood \quad
  Rich Wolski
}
\titlegraphic{\hfill\includegraphics[width=2.25cm]{ucsbseal_cmyk.pdf}}

\begin{document}
\maketitle

\metroset{numbering=none}
\begin{frame}<beamer>[noframenumbering]
  \frametitle{Outline}
  \tableofcontents
\end{frame}
\metroset{numbering=fraction}

\section{Motivation}
\begin{frame}{What is a System?}
Infrastructure software upon which applications are built
\\
\vspace{0.2in}
\pause

Operating Systems \note{process abstraction, named resource management} % lowest level of abstraction
   \begin{itemize}
       \item Process abstraction
       \item Multiplex physical hardware resources
       \item Partition and abstract \textbf{memory}
   \end{itemize}

\vspace{0.2in}
\pause
Embedded Systems, \pause Compilers, \pause Garbage Collectors, \pause Device Drivers, \pause File Systems
\end{frame}

\note[itemize]{
  \item example: compilers, garbage collectors, file systems, drivers, etc.
  \item need to be able to do this
  \item isolation third item for os
  \item files: hide peculiarites of disks/I/O decives, abstract model of device-indepdent files
  \item address space: virtual memory + protection
  \item process: running program container (registers, files, alarms, address space, etc.)
  \item they're all stateful, imperative things
}

\begin{frame}{Systems Languages}
\begin{itemize}[<+->]
    \item Past systems languages: \tiny{ALGOL, PL/I, Fortran, BCPL/B, C, Mesa/Cedar, Pascal/Modula-2/Oberon, C++, ...}
    \item C: the de-facto standard
        \begin{itemize}
            \item Data structure representation control \note{how it looks}
            \item Memory management control \note{where it is, how long it lives}
            \item Complete mutability via pointers \note{with systems without protected address space, can literally write anything; embedded have memory-mapped I/O}
            \item Performant
            \item Legacy \note{alas: have to deal with it}
        \end{itemize}
    \item C: the unsafe standard
        \begin{itemize}
            \item Unchecked array operations \note{buffer overflows}
            \item Pointers $\equiv$ arrays \note{hazardous pointer arithmetic}
            \item Unsafe casts \note{read/write arbitrary addresses, untyped}
            \item Aliasing \note{dangling pointers, double frees, null dereferences}
            \item \alert{Undefined behavior}
        \end{itemize}
\end{itemize}
\end{frame}

\note[itemize]{
\item ALGOL: formally defined syntax (OS: Burroughs, VME, Cambridge CAP)
\item PL/I: (OS: Multics)
\item Mesa, Cedar: rich exceptions, GC
\item Pascal (1971), Modula-2 (1982), Oberon (1988)
\item -- structured programming, records, pointers, dynamic allocation, information hiding, objects
% see mae3.org, mae4.org
}

\begin{frame}{Memory Safety}
\begin{itemize}[<+->]
    %\item Memory safety error: ``Any dereference of a pointer or subscripted array reference which reads or writes storage outside of the referent'' \footcite{austin_efficient_1994}
    \item Memory safety error: reads or writes outside the referent's storage
        \begin{itemize}
            \item Spatial: outside referent's \alert{address bounds}
            \item Temporal: outside referent's \alert{lifetime}
        \end{itemize}
    \item Ideal technique is
        \begin{itemize}
            \item Efficient and expressive % (hallmarks of C)
            \item Purely static % (no runtime overhead, we will not have if we want computer help)
            % note: checking buffer overflows is generally undecidable
            \item Precise % (not overly conservative)
            \item Automatic % (legacy code $\Rightarrow$ no source or memory layout change)
        \end{itemize}
%    \item Non-goals: secrecy, security, concurrency, type safety % but are natural consequences of safety
    \item Memory errors become \alert{type errors}, management happens at \alert{compile-time}
\end{itemize}
\end{frame}

\note[itemize]{
\item some of those goals from nagakatte 201*
\item undecidability of checking certain dynamic errors
\item hard to verify/prove invariants about b/c
\item -- casts + pointers make c essentially untyped (rondon: type system is only to help know size of bytes to read/write)
\item -- aliasing
\item make bad hard, useful easy
}

\AtBeginSection[]
{
  \metroset{numbering=none}
  \begin{frame}<beamer>[noframenumbering]
    \frametitle{Outline}
    \tableofcontents[currentsection]
  \end{frame}
  \metroset{numbering=fraction}
}

\AtBeginSubsection[]
{
  \metroset{numbering=none}
  \begin{frame}<beamer>[noframenumbering] 
    \frametitle{Outline}   
    \tableofcontents[currentsection,currentsubsection]
  \end{frame}
  \metroset{numbering=fraction}
}

\only<handout>{
  \addtocounter{framenumber}{1}
}

\section{Spatial Safety}

\begin{frame}[fragile]{Spatial Safety}
\begin{columns}[T]
\begin{column}{0.48\textwidth}
Always access within object's \alert{bounds}
\\
\vspace{0.2in}
\onslide<6->{
Some approaches:
\begin{itemize}
      \item Fat Pointers and Shadow Structures
      \item Referent Objects
      \item Dependent Types
%      \item Hardware
  \end{itemize}
}
\end{column}
\pause
\begin{column}{0.50\textwidth}
    \vspace{-0.2in}
\begin{lstlisting}[
    language=C,mathescape,basicstyle={\small\ttfamily}]
struct { int *arr; int i; } s;
s.arr = calloc(3, sizeof(int));
\end{lstlisting}

    \pause

    \begin{tikzpicture}
      \node[draw, fill=black!10!white, minimum width=12mm, minimum height=5mm, label=left:0x18] (n0) {0x40};
      \node[draw, fill=black!10!white, minimum width=12mm, minimum height=5mm, below=-\pgflinewidth of n0, label=left:0x1c] (n1) {.};
      \node[draw, minimum width=12mm, minimum height=5mm, below=-\pgflinewidth of n1,, label=left:0x20] (n2) {.};
      \node[draw, minimum width=12mm, minimum height=5mm, below=-\pgflinewidth of n2,, label=left:0x24] (n3) {.};
      \node[minimum width=12mm, minimum height=5mm, below=-\pgflinewidth of n3,,] (n4) {...};
      \node[draw, minimum width=12mm, minimum height=5mm, below=-\pgflinewidth of n4,, label=left:0x3c] (n5) {.};
      \node[draw, fill=black!30!white, minimum width=12mm, minimum height=5mm, below=-\pgflinewidth of n5,, label=left:0x40] (n6) {0};
      \node[draw, fill=black!30!white, minimum width=12mm, minimum height=5mm, below=-\pgflinewidth of n6,, label=left:0x44] (n7) {0};
      \node[draw, fill=black!30!white, minimum width=12mm, minimum height=5mm, below=-\pgflinewidth of n7,, label=left:0x48] (n8) {0};
      \node[draw, minimum width=12mm, minimum height=5mm, below=-\pgflinewidth of n8,, label=left:0x50] (n9) {.};
      \node[left = 1.5cm of n0] () {\textbf{s.arr}};
      \node[left = 1.5cm of n1] () {\textbf{s.i}};
      \onslide<3> \draw [->,black] (n0.east) to [out=315,in=45] (n6.east);
      \onslide<4> \draw[->,black] (n0.east) to [out=315,in=45] (n5.east);
      \onslide<5,6-> \draw[->,black] (n0.east) to [out=315,in=45] (n9.east);
    \end{tikzpicture}
\end{column}
\end{columns}

\end{frame}

\begin{frame}[fragile]{Spatial Safety Example}
\begin{columns}[T]
\begin{column}{0.52\textwidth}
\begin{lstlisting}[
    language=C,numbers=left,mathescape,basicstyle={\small\ttfamily},
    moredelim={**[is][{\btHL<1>}]{@1}{@}},
    moredelim={**[is][{\btHL<2>}]{@2}{@}},
    moredelim={**[is][{\btHL<3>}]{@3}{@}},
    moredelim={**[is][{\btHL<4>}]{@4}{@}},
    moredelim={**[is][{\btHL<5>}]{@5}{@}},
    moredelim={**[is][{\btHL<6>}]{@6}{@}},
    moredelim={**[is][{\btHL<7>}]{@7}{@}},
]
@2int find_token(int *data,@
@2               int *end,@
@2               int token)@
{
    @3int *p = data;@
    @4while (p < end) {@
    @5    if (*p == token) break;@
    @6    p++;@
    }
    @7return (*p == token);@
}
\end{lstlisting}
\end{column} 
\begin{column}{0.44\textwidth}

\onslide<8->
\footnotesize{
Potential pointer dereference problems:
  \begin{itemize}[<+->]
      \item Null
      \item Uninitialized
      \item Out-of-bounds \note{related: data might not be an array, unrelated}
      \item Manufactured \note{not shown here exactly}
  \end{itemize}
}
\end{column}
\end{columns}

\end{frame}

\subsection{Fat Pointers and Shadow Structures}

\begin{frame}[fragile]{Fat Pointers}
%Fat Pointers
\vspace{-0.1in}
\pause
\begin{columns}[T]
\begin{column}{0.64\textwidth}
\begin{itemize}[<+->]
 \item Added base and bound addresses \note{multiword, and other information}
 \item Instrument all pointers and functions
 \item Insert runtime checks \alert{before dereferences} \note{as opposed to during arith, let arithmetic happen as usual}
 \item Advantages:
     \begin{itemize}
        \item Quickly find and retrieve metadata \note{no costly db lookup}
        \item Complete spatial safety \note{only way to enforce complete spatial safety is tracking ptr bnds}
        \item No sub-object overflows \note{will explain next section}
     \end{itemize}
 \item Disadvantages:
     \begin{itemize}
        \item Breaks binary compatibility \note{calling convention of external/uninstrumented libraries, storing pointer in register, casting ptr to int and back}
        \item Metadata propagation \note{on pointer assignment}
        \item Code bloat, memory usage, runtime overhead \note{up to 200, because of added checks, bigger ptrs}
        \item Unsafe casts overwriting metadata
     \end{itemize}
\end{itemize}
\end{column}

\begin{column}{0.47\textwidth}

   \pause

\begin{lstlisting}[
    language=C,mathescape,basicstyle={\scriptsize\ttfamily},
    moredelim={**[is][\btHL<15>]{@1}{@}},
]
struct fptr { int *curr;
     int *base; int *bound; };
struct { @1struct fptr arr;@ int i; } s;
s.arr.curr = calloc(3, sizeof(int));
\end{lstlisting}

    \pause

    \begin{tikzpicture}
      \node[draw, fill=black!10!white, minimum width=12mm, minimum height=5mm, label=left:0x18] (n0) {0x40};
      \node[draw, fill=black!10!white, minimum width=12mm, minimum height=5mm, below=-\pgflinewidth of n0, label=left:0x1c] (n1) {0x40};
      \node[draw, fill=black!10!white, minimum width=12mm, minimum height=5mm, below=-\pgflinewidth of n1,, label=left:0x20] (n2) {0x4c};
      \node[draw, fill=black!10!white, minimum width=12mm, minimum height=5mm, below=-\pgflinewidth of n2,, label=left:0x24] (n3) {.};
      \node[minimum width=12mm, minimum height=5mm, below=-\pgflinewidth of n3,,] (n4) {...};
      \node[draw, minimum width=12mm, minimum height=5mm, below=-\pgflinewidth of n4,, label=left:0x3c] (n5) {.};
      \node[draw, fill=black!30!white, minimum width=12mm, minimum height=5mm, below=-\pgflinewidth of n5,, label=left:0x40] (n6) {0};
      \node[draw, fill=black!30!white, minimum width=12mm, minimum height=5mm, below=-\pgflinewidth of n6,, label=left:0x44] (n7) {0};
      \node[draw, fill=black!30!white, minimum width=12mm, minimum height=5mm, below=-\pgflinewidth of n7,, label=left:0x48] (n8) {0};
      \node[draw, minimum width=12mm, minimum height=5mm, below=-\pgflinewidth of n8,, label=left:0x4c] (n9) {.};
      \node[left = 1.0cm of n0] () {\textbf{s.arr.curr}};
      \node[left = 1.0cm of n1] () {\textbf{s.arr.base}};
      \node[left = 1.0cm of n2] () {\textbf{s.arr.bounds}};
      \node[left = 1.0cm of n3] () {\textbf{s.i}};
      \draw [->,black] (n0.east) to [out=315,in=45] (n6.east);
      \onslide<17-> \draw[->,black] (n1.east) to [out=315,in=45] (n6.east);
      \onslide<18-> \draw[->,black] (n2.east) to [out=315,in=45] (n9.east);
    \end{tikzpicture}
\end{column}
\end{columns}

\end{frame}
% write wrappers to remove inline metadata before calling outside function
% have to instrument malloc...

\note[itemize]{
\item Source-code annotations, for some
\item Impractical for large code bases
\item Two distinct pointers can point to same object and have diff base/bounds associated with them
\item Overcomes sub-object overflows b/c can be detected b/c uniqueness of each ptr
}

\begin{frame}{Fat Pointer Approaches}
SafeC \footcite{austin_efficient_1994}:
\begin{itemize}
  \item \alert{Safe} pointers have value, base, and size
  \item Complete spatial safety, if \alert{transparent} storage management and no safe pointer attribute \alert{manipulation} \note{no domain-specific allocator at user-level}
  \item 275\% space overhead, 2-6x runtime overhead, 0.35-3x code size overhead \note{no register allocated safe ptrs, compiler optimizations fail with additional checks}
  \item Some static optimization based on still-valid previous checks
\end{itemize}

\pause 

Cyclone \footcite{jim_cyclone:_2002}: \note{offshoot of TAL/popcorn}
\begin{itemize}
    \item Annotations for non-array vs array pointers (can specify size) % (never-null, fat, unitialized warning via control-flow analysis)
    \item Tagged unions and automatic tag injection \note{(e.g. instead of var-args, know exactly what's stored in thing callee passed $\texttt{printf}$)}
    %\item Can specify number of elements pointer points to \note{deputy is more general}
    %\item Arrays and strings converted automatically to fat
    %\item Statically validate and remove non-array fat pointers
    %\item 40\% runtime overhead
    %\item Uses regions + automatic memory management for temporal safety (free is a no-op) (see nice example)
    %\item never null don't need checks, use @; push back null checks from uses to their sources
    %\item Restrict arithmetic on regular pointers
    %\item parametric polymorphism, subtyping, static analysis to check for safety, adding run-time checks
\end{itemize}
\vspace{0.1in}

\end{frame}
\note[itemize]{
\item Cyclone requires more annotations by user than CCured
}

\begin{frame}{Fat Pointer Approaches}
 CCured \footcite{necula_ccured:_2002}
\begin{itemize}
    \item Separate pointers on usage (SAFE, SEQ, WILD) \note{safe: used w/o casts or ptr arith, either null or valid; seq: arith but no casts; wild: everything else}
    \item Whole-program inference to find as many SAFE then SEQ pointers as possible \note{via constraint collection}
    \item Reduce WILD pointers \footcite{necula_ccured:_2005} using physical subtyping \footcite{siff_coping_1999} for upcasts
    \item Special pointer RTTI carrying runtime type for downcasts \note{i.e. void* to anything, rttiOf mapping a type to node in physical subtyping hierarchy, store type info with ptr not object so they can point to inside of allocation unit}

    \note{wild: all regular ptrs can be treated as WILD (arbitrary casts, expensive dynamic checks)}
    \note{WILD/DYNAMIC: ptr type does not carry wit it the type of values pointed to}
    \note{seq: like fat pointers}
    \note{safe: if not null, guaranteed to point to memory are containing compatibly-typed object}
    \note{Relies on a garbage collector}
    \note{3\%-87\% runtime overhead increase}
    \note{quarantine wild/type-unsafe part from type-safe part}
\end{itemize}

% possible 2005: fig 5-6 

\pause

Fail-Safe C \footcite{oiwa_implementation_2009}:
\begin{itemize}
    \item Combines fat pointers w/ fat integers and virtual structure offsets \note{handles casts fine}
\end{itemize}

\vspace{0.2in}
\end{frame}

\note[itemize]{
\item CCured does translation of C programs to CCured most automatically
}

% use CCured syntax (2002, 2005 has a good translation definition)
\begin{frame}[fragile]{Fat Pointers Preventing Spatial Errors}
    \begin{columns}[T]
      \begin{column}{0.42\textwidth}
\begin{lstlisting}[language=C,numbers=left,mathescape,basicstyle={\scriptsize\ttfamily}]
int find_token(int *data,
               int *end,
               int token)
{
    int *p = data;
    while (p < end) {



        if (*p == token) break;


        p++;
    }

    return (*p == token);
}
\end{lstlisting}
\end{column}

\pause

\begin{column}{0.56\textwidth}

\begin{lstlisting}[language=C,numbers=left,mathescape,basicstyle={\scriptsize\ttfamily},
    moredelim={**[is][\btHL<3>]{@3}{@}},
    moredelim={**[is][{\btHL<4>}]{@4}{@}},
    moredelim={**[is][{\btHL<5>}]{@5}{@}},
    moredelim={**[is][{\btHL<6>}]{@6}{@}},
    moredelim={**[is][{\btHL<7>}]{@7}{@}},
    moredelim={**[is][{\btHL<8>}]{@8}{@}}
]
int find_token(int *@3SEQ@ data,
               int *@4SAFE@ end,
               int token)
{
    int *@3SEQ@ p = data;
    @5while (p.cur < end)@ {
        @6assert(p.base != 0 &&@
        @6         0 <= p.cur &&@
        @6         p.cur < p.bound);@
        @6if (*p.cur == token) break;@
        @7p.cur = p.cur + (1 * sizeof(int));@
        @7p.base  = p.base;@   // optimized out
        @7p.bound = p.bound;@  // "           "
    }
    @6...(repeat lines 7-9)...@
    @5return (*p.cur == token);@
}
\end{lstlisting}
\end{column}

\end{columns}
\end{frame}

\begin{frame}{Pointer-Based -- Shadow Structures}
MSCC \footcite{xu_efficient_2004}
\pause
\begin{itemize}
    \item \emph{Split} metadata from pointer, preserving layout \note{some libraries access based off of struct offset, or assume size of integer and ptr is same}
    \item Every value has linked shadow structure mirroring entire data structure
    \item Transform every function call to take additional metadata parameters
    \item Wrappers for external functions; cannot detect memory errors \note{Out-of-sync metadata after external library code}
  % \item Generate metadata for library routines returning newly allocated object
  % \item Does not handle arbitrary casts
  % \item Optimized encoding loses ability to detect sub-object overflows % no need to mention since we haven't talked about those flows yet
\end{itemize}
\end{frame}

% MSCC style
\begin{frame}[fragile]{Shadow Structures Example}
\begin{columns}
\begin{column}{0.42\textwidth}
\begin{lstlisting}[language=C,numbers=left,mathescape,basicstyle={\scriptsize\ttfamily}]




int find_token(
  int *data,
  int *end,
  int token)
{
    int *p = data;

    while (p < end) {

        if (*p == token) break;
        p++;
    }

    return (*p == token);
}
\end{lstlisting}
\end{column}

\pause

\begin{column}{0.56\textwidth}

\begin{lstlisting}[language=C,numbers=left,mathescape,basicstyle={\scriptsize\ttfamily},
    moredelim={**[is][{\btHL<3>}]{@3}{@}},
    moredelim={**[is][{\btHL<4>}]{@4}{@}},
    moredelim={**[is][{\btHL<5>}]{@5}{@}},
    moredelim={**[is][{\btHL<6>}]{@6}{@}},
]
@3struct ptr_info {@
@3    void *base;@
@3    unsigned long bound;@
@3};@
int find_token(
  int *data, @4ptr_info *data_info,@
  int *end,  @4ptr_info *end_info,@
  int token)
{
    int *p = data;
    @5ptr_info  p_info = *data_info;@
    while (p < end) {
        @6CHECK_SPATIAL(p, sizeof(*p), p_info);@
        if (*p == token) break;
        p++;
    }
    @6CHECK_SPATIAL(p, sizeof(*p), p_info);@
    return (*p == token);
}
\end{lstlisting}
    \end{column}
  \end{columns}
\end{frame}

\subsection{Referent Objects}

\begin{frame}[fragile]{Referent Objects}
    Objects \footcite{jones_backwards-compatible_1997}\footcite{ruwase_practical_2004}
\footnotesize
\begin{columns}[T]
\begin{column}{0.45\textwidth}
\begin{itemize}[<+->]
    \item Metadata about \alert{objects}, not pointers
    \item Global database relates every allocated address to corresponding object metadata \note{updating on alloc, dealloc, on function entry/exit for stack-based objects, program startup for global object}
    \item Every pointer to same object shares same metadata
    \item Bounds check on \alert{pointer arithmetic} \note{okay to have OOB ptr, if used for comparison, to ensure validity of derived pointers, point to same object}
    \item 2-12x overhead
\end{itemize}
\end{column}

\pause

\begin{column}{0.45\textwidth}
Advantages:
\begin{itemize}
  \item Compatible with uninstrumented code \note{no memory layout change; if pointer updated by untrusted module, metadata out-of-sync}
\end{itemize}
\pause
Disadvantages:
\begin{itemize}
  \item Special mechanisms to handle legal OOB pointers \note{since going 1 past array undereferenced is legal C}
  \item Splay-tree object lookup overhead
  \item Incomplete spatial safety: \alert{sub-object overflows} \note{type casts, mem corruption in objects, C allows ptr arith in struct fields}
\end{itemize}
\end{column}
\end{columns}

\vspace{0.1in}

\end{frame}

\note[itemize]{
\item jones and kelly: pad objects w/ extra byte to help with legal out-of-bounds ptr: ptr to a[N] is distinguishable from next object
\item check ptr arith and use, b/c result cannot refer to diff object from one from which it is originally derived
\item object for which the ptr is valid is only determined by checking ptr itself, looking it up in object table
\item Incomplete because cannot pad parameters, would change layout
\item cred: more generic solution
}

\begin{frame}[fragile]{The Sub-Object Problem}
\begin{lstlisting}[
    language=C,numbers=left,mathescape,basicstyle={\normalsize\ttfamily},
    moredelim={**[is][{\btHL<2>}]{@2}{@}},
    moredelim={**[is][{\btHL<3>}]{@3}{@}},
    moredelim={**[is][{\btHL<4>}]{@4}{@}},
    moredelim={**[is][{\btHL<5>}]{@5}{@}}
]
@2struct node {char str[3]; void (*func)(); };@
@3struct node *n = (struct node *) malloc(sizeof(node));@
@4char *s = n.str;@
@5strcpy(s, "bad!");@
\end{lstlisting}


\onslide<6->
    \begin{itemize}
        \item $\texttt{n}$ and $\texttt{s}$ have the same address $\Rightarrow$ map to same object in global database
        \item $\texttt{strcpy}$ will see $\texttt{s}$'s size as that of $\texttt{n}$
    \end{itemize}

\end{frame}

\note[itemize]{
\item Because nested objects (e.g., an array of structures) are assigned base and bound addresses
\item that span the entire allocated region, it is impossible to detect sub-object overflows if an
\item out-of-bounds pointer to an inner object remains within bounds of the outer object.
}

\begin{frame}[fragile]{Referent Objects Approaches}
SafeCode \footcite{dhurjati_backwards-compatible_2006}
    \begin{itemize}
        \item Use automatic pool allocation (APA) \footcite{lattner_automatic_2005}
        \item Use separate, \alert{smaller} data structures to store bounds metadata for \alert{each partition}
        \item 1.2x overhead
    \end{itemize}

\pause

Baggy Bounds Checking (BBC) \footcite{akritidis_baggy_2009}
    \begin{itemize}
    \item Compact bounds representation and efficient way to look up object bounds
    \item Align base addresses to be multiple of padded size
    \item Replace splay tree with small lookup table
    \item 0.6x overhead \note{on SPECINT 2000 benchmark}
    \end{itemize}
\vspace{0.2in}

\end{frame}
\note[itemize]{
   \item builds on JK, APA: partitions memory based on static points-to analysis
   \item Trades memory for performance, fastest object bounds checker
   \item partition memory into slots with slot_size bytes (table has entry per slot rather than per byte)
   \item Pad every object s.t. size is power of two, align base addresses to be multiple of their padded size
   \item mark if OOB to prevent later dereference
   \item handle OOB withint slot_size/2 bytes from original object
}

\begin{frame}[fragile]{Referent Objects Example}
\footnotesize
\begin{columns}[T]
\begin{column}{0.42\textwidth}
\begin{lstlisting}[language=C,numbers=left,mathescape,basicstyle={\scriptsize\ttfamily}]
int find_token(int *data,
               int *end,
               int token)
{
    int *p = data;
    while (p < end) {
        if (*p == token) break;




        p++;
    }
    return (*p == token);
}
\end{lstlisting}
\end{column}

\pause

\begin{column}{0.56\textwidth}
\begin{lstlisting}[
    language=C,numbers=left,mathescape,basicstyle={\scriptsize\ttfamily},
    moredelim={**[is][{\btHL<3>}]{@3}{@}},
    moredelim={**[is][{\btHL<4>}]{@4}{@}},
    moredelim={**[is][{\btHL<5>}]{@5}{@}},
    moredelim={**[is][{\btHL<6>}]{@6}{@}},
]
int find_token(int *data,
               int *end,
               int token)
{
    int *p = data;
    while (p < end) {
        if (*p == token) break;
        @3int *q = p + 1;@
        @4int size = 1 << TABLE[p>>4];@
        @5int base = p & ~(size - 1);@
        @6assert(q >= base && q - base < size);@
        p++;
    }
    return (*p == token);
}
\end{lstlisting}
%Lines 4-6 are equivalent to $\texttt{assert((p^q)>>TABLE[p>>SLOT_SIZE] == 0)}$.
\end{column}

\end{columns}
\end{frame}

\begin{frame}{Softbound}
%\footnotesize
Softbound \footcite{nagarakatte_softbound:_2009}
    \begin{itemize}[<+->]
        \item Base and bound metadata for each pointer, stored in disjoint metadata table
        \item Total spatial safety of \alert{pointer-based} approaches
        \item Source compatibility, separate compilation of \alert{object-based} approaches
        \item Runtime bounds checks on each dereference
        \item Propagate metadata as extra arguments
        \item Arbitrary casts allowed
        \item 67\% overhead
    \end{itemize}
\end{frame}
\note[itemize]{
    \item Must protect every module for *complete* safety (untrusted module can modify ptr, table is out-of-sync)
    \item Only store metadata in disjoint space for pointers in memory, perform metadata loads/stores only when loaded/stored value from/into memory loc is a ptr
    \item Wrappers for commonly used libraries
    \item like pointer-based and object approaches, two-level trie data structure in disjoint metadata space
    \item eliminate sub-object problem
    \item just check ptr involved in casts on derfernce w.r.t to its metadata, no way to overwrite metadata
    \item to function calls, wrappers for commonly-used libraries, must transform every func decl/call-site, or shadow stack (nagarakatte thesis)
}

\subsection{Dependent Types}

\begin{frame}{Dependent Types}
    Dependent types are \emph{typed-valued functions} \footcite{pierce_advanced_2005}
    % type family of vectors
    % dependent product type
\pause
\begin{center}
   $\texttt{Vector : Nat}\rightarrow\texttt{Type}\rightarrow\texttt{Type}$
   \pause
   \\
   $\texttt{nil : Vector 0 a}$
   \\
   \pause
   $\texttt{cons : }\Pi\texttt{n:Nat.a}\rightarrow\texttt{Vector n a}\rightarrow\texttt{Vector (n+1) a}$
   \\
   \pause
   $\texttt{(cons 'a' (cons 'b' nil)) : Vector 2 Char}$
   \\
   \pause
   $\texttt{head : }\Pi\texttt{n:Nat.Vector (n+1) a -> a}$
   \\
   \pause
   $\texttt{head nil}$ \pause $\Rightarrow$ \textcolor{red}{Rejected!}
\end{center}
    \pause
    \begin{itemize}
        \item Based on type theory work by Martin-L{\"o}f \footcite{martin-lof_constructive_1984}
        \item \alert{Undecidability} of type checking: arbitrary computation to check type equality
        \item Work on defining equality and restricting forms of index terms
    \end{itemize}
    \vspace{0.1in}
\end{frame}
\note[itemize]{
  \item Declares a family of types
  \item We say that Vect is indexed over Nat and parameterised by Type
}

\begin{frame}{Early Uses of Dependent Types}
Dependent ML \footcite{xi_eliminating_1998} and Cayenne \footcite{augustsson_cayennelanguage_1998}
\begin{itemize}
    \item Reduce static array bound checking to constraint satisfiability
    \item DML uses \emph{indexed types}: limit indices to linear integer and boolean expressions; compile-time decidable
    \item Cayenne has \emph{no restrictions} on types: undecidability of arbitrary expression equivalence and thus type checking \note{haskell-like}
    %\item User provides types of recursive procedures and inductive datatypes
\end{itemize}
Xanadu \footcite{xi_imperative_2000}
\begin{itemize}
    \item Imperative environment
    \item Restrict index expressions in types to integer constraint domain \note{like DML}
%    \item Defines what it means for a type to equal/coerce into another
%    \item Programmer must supply state type in order to type conditionals and loops
\end{itemize}
    \vspace{0.2in}
\end{frame}
\note[itemize] {
\item the goal here is just to show that FuncLangs did it first, but they aren't always practical b/c of undecidability, difficult to use
\item infuttered form, you get CoC
\item decidable logic
\item note{C-like}
}

% \item Allows trusted cast mechanism to supress errors

\begin{frame}{Dependent Types in Imperative Languages}
    SafeDrive \footcite{zhou_safedrive:_2006} and Deputy \footcite{condit_dependent_2007}\footcite{anderson_static_2007}
%\vspace{-0.1in}
    \begin{itemize}[<+->]
        \item User-added annotations relating pointers to bounds
            \begin{itemize}
                \item \texttt{safe}, \texttt{sentinel}, \texttt{count(n)}, \texttt{bound(lo,hi)}
                \item Use constants/variables/field names in immediately enclosing scope \note{n, lo, hi are expressions}
            \end{itemize}
        %\item allow programmers to specify common c programming idioms involving ptr arith and union types
        %\item compiler checks safety of most pointer and union operations
        \item Three-phase pass over annotated C programs, emits C code
            \begin{enumerate}
                \item Automatic addition of bounds annotations for pointer types \note{fresh local var to hold bounds}
                \item Flow-insensitive type checking (insert run-time checks; helps decidability) \note {hybrid means SA for type determination, run-time checks when SA precision is too low}
                \item Flow-sensitive check optimization
            \end{enumerate}
        \item More C Support with dependent union tags \note{c idiom: tag indicates union field in use}, Safe TinyOS \footcite{cooprider_efficient_2007} %\footcite{hill_system_2000}, \footcite{gay_nesc_2003}, \footcite{cooprider_efficient_2007}
    \end{itemize}
    \vspace{0.2in}
\end{frame}

\note[itemize]{
  \item they're really *hybrid* approach: try your best, then offer runtime checks
  \item track info required to place run-time checks (asserts) in program
  \item  flow- path-insensitive, bolvious to aliasing, heap updates, data structures
  \item local type inference
  \item static verifier discharges checks at compile-time
% about annotations
  \item safe = ptr is null/pts to single element of base type
  \item sentintel = only for comparisons
  \item count(n) = null or point to an array of at least nelements of base type
  \item bounds(lo,hi) = null or point into mem region w/ low address lo, hi is first element beyond region end
  \item: invariant is that ptrs null, or within stated bounds
% dependent union tags
  \item c idiom: tag indicates union field in use
  \item Specify the condition for each union field to be usable
  \item \emph{$\text{union}_n\ \tau_1\ ...\ \tau_n\ e_1\ ...\ e_n$}
  \item $x:\text{struct }\{ tag:\text{int}\;\ u: \text{union}_2 \text{ int (ref int)} \ (\textit{tag} \geq 2)\ (\textit{tag} = 1) \}$
% generalized ptr
  \item Generalize array constructor to possibly-null bounded pointer
  \item C-style pointer arithmetic via $\oplus$ operator
  \item $x:\text{ptr int } x\ (x \oplus n)$ $\texttt{// n integer area starting at x}$
  \item $x:\text{ptr int } x\ e\ $ $\texttt{// from x to e}$
}

% go through an example to get a feeling, then after explaining judgments, explain it again stepping through the judgments
\begin{frame}[fragile]{Deputy Example for Spatial Safety}
\begin{columns}[T]
\begin{column}{0.40\textwidth}
\begin{lstlisting}[language=C,numbers=left,mathescape,basicstyle={\scriptsize\ttfamily}]
int find_token(int *data,
               int *end,
               int token)
{


    int *p = data;
    while (p < end) {


        if (*p == token) break;

        p++;
    }


    return (*p == token);
}
\end{lstlisting}
\end{column}

%\scriptsize{$\colorbox{blue!30}{Code added during automatic dependency inference}$}

\pause

\begin{column}{0.58\textwidth}
\begin{lstlisting}[
    language=C,numbers=left,mathescape,basicstyle={\scriptsize\ttfamily},
    moredelim={**[is][{\btHL[onslide=<3->{fill=blue!50}]}]{@3}{@}},
    moredelim={**[is][{\btHL[onslide=<4->{fill=blue!50}]}]{@4}{@}},
    moredelim={**[is][{\btHL[onslide=<5->{fill=blue!50}]}]{@5}{@}},
    moredelim={**[is][{\btHL[onslide=<6->{fill=red!50}]}]{@6}{@}},
    moredelim={**[is][{\btHL[onslide=<7->{fill=red!50}]}]{@7}{@}},
    moredelim={**[is][{\btHL[onslide=<8->{fill=red!50}]}]{@8}{@}},
    moredelim={**[is][{\btHL[onslide=<9->{fill=red!50}]}]{@9}{@}},
]
int find_token(int * @3bound(data, end)@ data,
               int * @4sentinel@ end,
               int token)
{
    @6assert(data != NULL);@
    @6assert(end != NULL);@
    int *p @5bound(p, end)@ = data;
    while (p < end) {
        @7assert(p != NULL);@
        @7assert(p < end);@
        if (*p == token) break;
        @8assert(p <= p + 1 <= end);@
        p++;
    }
    @9assert(p != NULL);@
    @9assert(p < end);@
    return (*p == token);
}
\end{lstlisting}
\end{column}
\note[itemize]{
\item Let's assume that the code is fully annotated (automatic generation coming later)
\item No arithmetic/comparisons on null ptrs
\item 8 assert: legal to point past end
\item 9 illegal to access
}

\end{columns}
\end{frame}
% \note{All checks here would be eliminated}

% \note{here I need to exaplin what a grammer is for them}
\begin{frame}{Abstract Syntax, For Your Consideration}
\begin{gather*}
    x,y \in \text{Variables}
    \quad
    \text{op} \in \text{Binary ops}
    \quad
    n \in \text{Integers}
    \quad
    \text{comp} \in \text{Comparison Ops}
\end{gather*}

\vspace{0.2in}

\begin{columns}[T]
\begin{column}{0.45\textwidth}
\begin{alignat*}{2}
\textit{Ctors}\quad &C &&::= \text{int} \altm \text{ref} \altm \text{array}
\\
\textit{Types}\quad &\tau &&::= C \altm \tau_1 \ \tau_2 \altm \colorIt{blue}{n1}{\tau\ e}
%\\
%\text{Kinds}\quad & \kappa &&::= \text{type} \altm \text{type} \rightarrow \kappa \altm \tau \rightarrow \kappa
\\
\textit{L-exprs} \quad &l &&::= x \altm *e
\end{alignat*}
\end{column}

\pause

\begin{column}{0.45\textwidth}
\begin{alignat*}{2}
\textit{Exprs}\quad &e &&::= n \altm l \altm e_1 \text{ op } e_2
\\
\textit{Cmds}\quad &c &&::= l := e \altm \colorIt{red}{n2}{\text{assert}(\gamma)} \altm c_1;c_2 \altm ...
\\
%\textit{Cmds}\quad &c &&::= \text{skip} \altm c_1;c_2 \altm l := e \altm \colorIt{red}{n2}{\text{assert}(\gamma)} \altm
%\\
%~ & ~ && \text{let } x : \tau = e \text { in } c \altm \text{let } x = \text{new } \tau(e) \text{ in } c
%\\
\colorIt{green}{n3}{\text{Preds}} \quad &\gamma &&::= e_1 \text{ comp } e_2 \altm \text{true} \altm \gamma_1 \wedge \gamma_2
\end{alignat*}
\end{column}
\end{columns}

\vspace{0.1in}
\end{frame}

%\scriptsize{$\colorbox{green!30}{Programmer-supplied annotation}$}

\note[itemize]{
\item kind of int is "type"
\item kind of ref is "type -> type"
\item loops, conditionals omitted for flow-insenstivei type system
\item Types restricted to only contain expressions using constatns, local variables, arbitrary arithmetic operators
\item Type environment is a predicate on the state of the program
}

\begin{frame}{Typing Rules, For Your Consideration}
\note{local expression: expressions formed using constants, local variables, arbitrary arithmetic operators (no memory deferences}
\note{expressions in types are restricted to local expressions}

Local Expressions: $\Gamma \vdash_{L} e : \tau$
%Local Expressions: $\colorIt{red}{n1}{\Gamma } \colorIt{blue}{n2}{\vdash_{L} e} \colorIt{orange}{n3}{: \tau}$
\vspace{-0.1in}
\pause
\begin{center}
\begin{tabular}{cc}
\infer[(\textsc{local\ name})]
{\colorIt{red}{n1}{\Gamma } \colorIt{blue}{n2}{\vdash_{L} x} \colorIt{orange}{n3}{: \tau }}
{\colorIt{green}{n4}{\Gamma(x) = \tau}}
    &
\infer[(\textsc{local\ num})]
{\Gamma \vdash_{L} n : int }
    {}
\end{tabular}
\end{center}

\pause

Non-local Expressions: $\Gamma \vdash e : \tau \Rightarrow \gamma$
%Non-local Expressions: $\colorIt{red}{n1}{\Gamma } \colorIt{blue}{n2}{\vdash e } \colorIt{orange}{n3}{: \tau \Rightarrow \gamma}$
\pause
\vspace{-0.1in}
\begin{center}
\begin{tabular}{c}
\infer[(\textsc{deref})]
{\colorIt{red}{n1}{\Gamma} \colorIt{blue}{n2}{\vdash *e } \colorIt{orange}{n3}{: \tau \Rightarrow \gamma}}
{\colorIt{green}{n4}{\colorIt{red}{n5}{\Gamma} \colorIt{blue}{n6}{\vdash e} \colorIt{orange}{n7}{: \text{ref } \tau \Rightarrow \gamma}}}
\end{tabular}
\end{center}

\pause

Commands: $\Gamma \vdash c \Rightarrow c'$
%Commands: $\colorIt{red}{n1}{\Gamma } \colorIt{blue}{n2}{\vdash c } \colorIt{orange}{n3}{\Rightarrow c'}$
\vspace{-0.1in}
\pause
\begin{center}
\begin{tabular}{c}
\infer[(\textsc{var\ write})]
{\colorIt{red}{n1}{\Gamma } \colorIt{blue}{n2}{\vdash x \coloneqq e } \colorIt{orange}{n3}{\Rightarrow \text{assert}(\bigwedge_{y \in \text{Dom}(\Gamma)}\gamma_y);\ x \coloneqq e }}
{\colorIt{green}{n4}{x \in \text{Dom}(\Gamma)} \qquad
  \colorIt{green}{n8}{\forall(y:\tau_y) \in \Gamma,\ \colorIt{red}{n5}{\Gamma } \colorIt{blue}{n6}{\vdash y[e/x]} \colorIt{orange}{n7}{:\tau_y[e/x] \Rightarrow \gamma_y }}}
\end{tabular}
\end{center}

\pause
\end{frame}

\begin{frame}{The Interesting Rules}
Dereferencing
\begin{center}
\begin{tabular}{c}
\infer[(\textsc{array\ deref})]
{\Gamma \vdash *e; \tau \Rightarrow \gamma_e \wedge (0 < e_{len})}
{\Gamma \vdash e: \text{array } \tau \ e_{len} \Rightarrow \gamma_e}
\end{tabular}
\end{center}

\vspace{0.1in}
\pause

Arithmetic
\begin{center}
\begin{tabular}{c}
\infer[(\textsc{array\ arith})]
{\colorIt{red}{n1}{\Gamma} \colorIt{blue}{n2}{\vdash e + e'} \colorIt{orange}{n3}{: \text{array } \tau\ (e_{len} - e') \Rightarrow \gamma_e \wedge \gamma_e' \wedge (0 \leq e' \leq e_{len})}}
{\colorIt{green}{n7}{\colorIt{red}{n4}{\Gamma} \colorIt{blue}{n5}{\vdash e} \colorIt{orange}{n6}{: \text{array } \tau\ e_{len} \Rightarrow \gamma_e}}
 \quad \colorIt{green}{n8}{\colorIt{red}{n9}{\Gamma} \colorIt{blue}{n10}{\vdash e'} \colorIt{orange}{n11}{:\text{int} \Rightarrow \gamma_{e'}}}}
\end{tabular} 
\end{center}

\end{frame}
% var-write is a key contribution of type system
% $\textsc{var\ write}$ rule responsible for updates to variables in presence of dependent type variables, verifying assignment does not break any dependencies in current scope
% For example, we can add rules for arrays:
% also permits a coercion judgement for allowing an array to be used where a smaller one is expected

\begin{frame}{Dependent Types in Imperative Languages}

T\'yr \footcite{de_araujo_tyr:_2016}
\vspace{-0.1in}
\begin{itemize}
    \item Augments LLVM IR with dependent pointer types
    \item Uses programmer annotations insert run-time bounds checks
    \item LLVM optimizations remove always-true checks; error if always-false
\end{itemize}

\pause

Checked C \footcite{ruef_checked_2017}
\vspace{-0.1in}
    \begin{itemize}
%        \item Focus on interoperablity, usability, performance
%        \item Discovers safely-used ptrs, rewriting them to be checked
%        \item Normal C pointers
        \item Extend C with two \emph{checked pointer types}: $\texttt{\_Ptr<T>}$ and $\texttt{\_Array\_ptr<T>}$ \note{automatically rewrite code to use when possible; no change to pointers}
        \item Associated bounds expressions indicating where bounds are stored
        \item Isolate (un)safe code with \emph{checked code regions}  \note{prevent unchecked pointer usage and certain casts, at file/func/block level; cannot blame checked code for violations}
%        \item Cannot blame checked code for violation
%        \item Derefence-only (no arith.) pointer, arithmetic-supporting (possibly null-terminated) pointer with bounds in type
%        \item $\verb{_Ptr<T>}$, a pointer for dereference only (no arith)
%        \item $\verb{_Array_ptr<T>}$ and $\verb{_Nt_array_ptr<T>_}$, supporting arith w/ bounds declarations in type (latter is null-terminated)
%       \item Compiler confirms restrictions maintained, inserts checks on ptr access
    \end{itemize}

\pause

Low* \footcite{protzenko_verified_2017}
\vspace{-0.1in}
\begin{itemize}
    \item DSL for verified, efficient low-level programming in F* %(ML-like language with dependent types)
    \item Write F* syntax against library modelling lower-level view of C memory
\end{itemize}
\vspace{0.1in}
\end{frame}
\note[itemize]{
\item Translate Low* to CompCert Clight
\item Goal: write efficient \& verified C in a high-level language
\item Model arrays by implementation abstract buffer type using references by hyper-stacks
}

\begin{frame}{Quick Spatial Recap}
Spatial Safety
\begin{itemize}[<+->]
    \item Arrays and pointers \note{Unsafe languages like C conflate arrays and pointers, casts eliminate any semblance of type guarantees, types are just for size}
    %\item Undecidable \note{Completely detecting buffer overflow in general is undecidable, so we try runtime things}
    \item Fat pointers \note{complete spatial safety at cost of compatibility}
    \item Referent objects \note{partial safety and better compatibility}
    \item Dependent types \note{prove in-bounds access at compile-time, emit run-time checks for full safety, Solution: restrict C, add bounds, runtime checks}
\end{itemize}
\end{frame}

\section{Temporal Safety}

\begin{frame}[fragile]{Temporal Safety}
Prevent accessing object that has been \alert{previously deallocated} \note{could be allocated again, but your pointer is wrong}
\\
\pause

\begin{itemize}
   \item Capabilities and locks
   \item Effects and regions
   \item Linear types and ownership
\end{itemize}
\end{frame}

\begin{frame}[fragile]{Temporal Safety Example}
\begin{columns}[T]
\begin{column}{0.52\textwidth}
\begin{lstlisting}[
    language=C,numbers=left,mathescape,basicstyle={\scriptsize\ttfamily},
    moredelim={**[is][{\btHL<1>}]{@1}{@}},
    moredelim={**[is][{\btHL<2>}]{@2}{@}},
    moredelim={**[is][{\btHL<3>}]{@3}{@}},
    moredelim={**[is][{\btHL<4>}]{@4}{@}},
    moredelim={**[is][{\btHL<5>}]{@5}{@}},
    moredelim={**[is][{\btHL<6>}]{@6}{@}},
    moredelim={**[is][{\btHL<7>}]{@7}{@}},
    moredelim={**[is][{\btHL<8>}]{@8}{@}},
    moredelim={**[is][{\btHL<9>}]{@9}{@}},
    moredelim={**[is][{\btHL<10>}]{@10}{@}},
    moredelim={**[is][{\btHL<11>}]{@11}{@}},
    moredelim={**[is][{\btHL<12>}]{@12}{@}},
]
@1int attach(struct sock *sk) {@
@2    if (sk->bad) {@
@3        free(sk); return 1;@
    }
@4    return 0;@
}
@5void mq_notify(sigevent *n) {@
@6    struct sock_t *sock;@
@7    while (n->try) {@
@8        sock = malloc_sock(n->info);@
@9        if (attach(sock)){@
               //sock = NULL;
@10            break;@
        }
    }
@11    if (sock) free(sock);@
}
\end{lstlisting}
\end{column} 
\begin{column}{0.44\textwidth}

\footnotesize{
\onslide<13->
Potential pointer dereference problems:
  \begin{itemize}
      \item Double frees
      \item Dangling pointers
  \end{itemize}
\pause

\onslide<14->
A Real Bug
 \begin{itemize}
    \item Linux Kernel in \texttt{ipc/mqueue.c}
    \item July 2017
    \item \url{https://bugzilla.redhat.com/show_bug.cgi?id=1470659}
  \end{itemize}

\pause

\onslide<15->
Goals:
\begin{itemize}
    \item Good: Detecting use-after-free
    \item Better: Eliminating free entirely \note{so 'use-after' has no meaning}
\end{itemize}
}
\end{column}
\end{columns}

\end{frame}

\begin{frame}{A Comment on Garbage Collection}
Garbage collection
\pause
    \begin{itemize}[<+->]
        \item Relinquish control of object location and layout to runtime \note{disallow free}
        \item Complete temporal safety, but...
        \begin{itemize}
            \item Non-zero overhead
            \item Drag
            \item Loss of real-time guarantees/predictability
            \item Reduced reference locality, increased page fault/cache miss rates
        \end{itemize}
        %\item Non-zero overhead, sub-par memory use (drag) \note{time between unreachable and reclaimed}
        \item Some spatial approaches (e.g. Fail-Safe C, CCured) use Boehm-Demers-Weister \footcite{boehm_garbage_1988}
           \note{heapify local variables (prevent dangling ptrs to stack objects) so they are managed by GC}
        %\item Others like Cling \footcite{akritidis_cling:_2010} rely on a special allocators
        %    \note{Replace malloc: address space re-use only among objects of same type and alignment (type-safe memory re-use)}
    \end{itemize}
\end{frame}
% \footcite{austin_efficient_1994}
% \footcite{jones_backwards-compatible_1997}
% \footcite{xu_efficient_2004}

\subsection{Capabilities and Locks}

\begin{frame}{Capabilities and Locks}
    \pause
    \begin{itemize}
        \item SafeC, MSCC %, Patil 1997, Yong 2003
            \begin{itemize}
              \item Unique capability associated with each memory block
              \item Stored in capability store, marked invalid on free
              \item Check if pointer's capability copy is still valid on dereference
            \end{itemize}
        \pause
        \item CETS \footcite{nagarakatte_cets:_2010}\footcite{nagarakatte_everything_2015}
            \begin{itemize}
                \item Each allocation has unique (never reused) key and lock address \note{key associated with all pointers to object}
                \item Freeing allocated object changes value at lock location, so key and lock value don't match
            \end{itemize}
        \pause
        \item Memsafe \footcite{simpson_memsafe:_2013}
            \begin{itemize}
                \item Set bounds of deallocated pointer to invalid value % subsequent dereferences raise exception as bounds is invalid
            \end{itemize}
    \end{itemize}
\end{frame}

\begin{frame}[fragile]{Lock Example}

\begin{columns}[T]
\begin{column}{0.48\textwidth}
\begin{lstlisting}[
    language=C,numbers=left,mathescape,basicstyle={\tiny\ttfamily},
    moredelim={**[is][{\btHL<1>}]{@1}{@}},
    moredelim={**[is][{\btHL<2>}]{@2}{@}},
    moredelim={**[is][{\btHL<3>}]{@3}{@}},
]
int attach(struct sock *sk,
@1      key_t sk_key,@
@1      lock_t *sk_lock_addr) {@
@2    if (sk_key != *sk_lock_addr)@
@2      abort();@
    if (sk->bad) {
@3        if (Freeable_ptrs_map.lookup(sk_key) != sk)@
@3          abort();@
        free(sk);
@3        *sk_lock_addr = INVALID_KEY;@
@3        deallocate_lock(sk_lock_addr);@
        return 1;
    }
    return 0;
}
\end{lstlisting}
\end{column}
\begin{column}{0.50\textwidth}
\begin{lstlisting}[
    language=C,numbers=left,mathescape,basicstyle={\tiny\ttfamily},
    moredelim={**[is][{\btHL<4>}]{@1}{@}},
    moredelim={**[is][{\btHL<5>}]{@2}{@}},
    moredelim={**[is][{\btHL<6>}]{@3}{@}},
]
void mq_notify(sigevent *n) {
    struct sock_t *sock;
@1    key_t sock_key;@
@1    lock_t *sock_lock_addr;@
    while (n->try) {
        sock = malloc_sock(n->info);
@2        sock_key = Next_key++;@
@2        sock_lock_addr = allocate_lock();@
@2        *(sock_lock_addr) = sock_key;@
@2        Freeable_ptrs_map.insert(sock_key, sock);@
        if (attach(sock)){
            break;
        }
    }
    if (sock) {
@3        if (Freeable_ptrs_map.lookup(sock_key) != sock)@
@3          abort();@
        free(sock);
@3        *(sock_lock_addr) = INVALID_KEY;@
@3        deallocate_lock(sock_lock_addr);@
    }
}
\end{lstlisting}
\end{column}
\end{columns}
\end{frame}

\subsection{Effects and Regions}

\note[itemize]{
  \item Can we get the \alert{compiler} to do the memory deallocation for us? % will require a restriction of paradigms/cooperation by the programmer
  \item Can we statically determine the lifetimes of all values? \note{like automatic variables put on/popped off stack}
  \item let's begin with effects and regions
}

\begin{frame}{Effects}
    Effect types describe the \emph{effects} of the computation leading to a value \footcite{pierce_advanced_2005}
    \\
    \begin{itemize}
        \item Opening a file
        \item Modifying an object
    \end{itemize}
\end{frame}

\begin{frame}{Effects}
    Fluent Languages\footcite{gifford_integrating_1986}, MFX\footcite{lucassen_polymorphic_1988}
  \begin{columns}[T]
    \begin{column}{0.68\textwidth}
        \begin{itemize}[<+->]
            \item Mix functional and imperative languages \note{determine which can be parallelized}
            \item Every expression has a \alert{effect class}, restricting which subroutines or sublanguage features it may use
            \item Identify \alert{referentially transparent} expressions via a side effect lattice \note{memoizable, for concurrency}
            \item Effect masking: inference to delimit regions of memory and their lifetimes \note{lets imperative program fragments be embedded in functional programs}
            %\item side-effects that cannot be observed outside of a given expression can be masked by the effect system
            %\item effect masking lets imperative program fragments be embedded in functional programs
        \end{itemize}
    \end{column}

    \pause

    \begin{column}{0.30\textwidth}
    \vspace{-0.2in}
    \footnotesize{
    \begin{tikzpicture}
        \node (WRA) at (0,3) {$\{W\ R^{?}\ A^{?}\}$};
        \node (RA)  at (0,2) {$\{R\ A^{?}\}$};
        \node (A)   at (0,1) {$\{A\}$};
        \node (e)   at (0,0) {$\{\}$};
        \node[draw] at (2,3) {$\text{Procedure}$};
        \node[draw] at (1.95,2) {$\text{Observer}$};
        \node[draw] at (1.95,1) {$\text{Function}$};
        \node[draw, right= of e] {$\text{Pure}$};
      \draw (e) -- (A) -- (RA) -- (WRA);
    \end{tikzpicture}

    \vspace{0.1in}
    \pause
    E.g.
    \begin{itemize}
        \item \texttt{update} is Procedure
        \item \texttt{nth} is Observer
        \item \texttt{arrayCreate} is Function
        \item \texttt{length} is Pure
    \end{itemize}
    }
    \end{column}
  \end{columns}
  \vspace{-0.3in}
\end{frame}
\note[itemize]{
% fluent
\item observer: not cause side-effects; cannot perform assignments or call procedure
\item function: not cause side-effcts, value is not affected by side-effects; no assignments, cannot call up, cannot access update-able variables
\item pure: referentially transparent: does not cause side-ffects, value not affected by side-effects, returns same value each time evaluated
\item pure:  no asignments, no accessing updating variables, cannot allocate updateable variables
\item cannot call UP the lattice
% MFX:
\item show soundness of effect type system for FX language
\item use effect information to constraint parallel evaluation of expressions
\item effect-masking: side-effects taht cannot be observed outside of the given expression
\item show how type and effect info can be used to implement certain kinds of side-effect analyses
\item language extends work with:
\item effect classes
\item full abstraction over effects
\item unmaskable side-effects
\item soundness proof
}

\begin{frame}{Effects}
    Type and Effect Systems \footcite{nielson_type_1999}
    \begin{itemize}[<+->]
        \item Extend the simply-typed lambda calculus with annotations
        \item Various forms of analyses, incl. Side Effect Analysis and \alert{Region} Inference
    \end{itemize}

    \pause

\vspace{-0.2in}
\begin{alignat*}{2}
    \text{Exprs}\quad &e &&::= c \altm x \altm \texttt{fn}_{\pi}\ x\ \Rightarrow e \altm e_1\ e_2 \onslide<5->{\altm \texttt{new}_{\pi}\ x\ :=\ e_1\ \texttt{in}\ e_2} \onslide<6->{\altm!x} \onslide<7->{\altm x := e}
\\
    \text{Types}\quad &\tau &&::= \texttt{int} \altm \texttt{bool} \only<4-8>{\altm \tau_1 \rightarrow \tau_2} \only<9->{\altm \tau_1 \xrightarrow[]{\phi} \tau_2} \only<8>{\altm \tau\ \texttt{ref}} \only<9->{\altm \tau\ \texttt{ref}\ \varrho}
\onslide<10->{\\ \only<10-11>{\text{Effect}}\only<12->{\colorIt{blue}{n1}{\text{Region}}}\quad &\phi &&::= \{!\pi\} \altm \{\pi:=\} \altm \{\texttt{new}\ \pi\} \altm \phi_1 \cup \phi_2 \altm \emptyset}
\onslide<11->{\\ \text{PrgPts}\quad &\varrho &&::= \{\pi\} \altm \varrho_1 \cup \varrho_2 \altm \emptyset}
\end{alignat*}

\onslide<13->{
\begin{center}
    $\Gamma \vdash e\ :\ \tau \ \&\ \phi$
\end{center}
}

\vspace{0.1in}
\end{frame}
\note[itemize]{
\item region inference: compile-time determination of when to allocate regions, in which to allocate data
\item Ref. variable represented by set of points where it could have been created (\alert{region})
\item Annotate functions with actions they take (access, assignment, creation)
\item we call \pi the creation point
\item \tau ref \varrho is type of loc created at one of the program points in region \varrho
}

\begin{frame}{Region-Based Memory Managememt}
    Regions \footcite{tofte_region-based_1997}
    \begin{itemize}[<+->]
        \item Divide heap into stack of sub-heaps (i.e. regions) \note{unbounded subheaps}
        \item Regions grow on individual allocation; \alert{entire} region deallocated \note{no individual object deallocation}
        \item Explicit region polymorphism \note{regions can be given as arguments to functions at runtime}
        \item Sound type system \alert{guarantees safety} of deallocations \note{both alloc and dealloc are explicit}
        \item Region inference identifies:
            \begin{itemize}
                \item Points where entire regions are allocated and deallocated
                \item Into which region values should go
            \end{itemize}
        %\item Used in ML Kit
        \item Unreasonable object lifetimes due to LIFO ordering of region lifetimes %(lifetime of expression must coincide with time takes to execute some subexpression)
        %\item new regions must be removed before older ones (but you can put a value in any region)
    \end{itemize}
    \vspace{0.1in}
\end{frame}

% Goal: realistic region support for C, but don't give tons of safety guarantees: more to help
%       b/c they allow arbitrary region intermingling, they need reference counts
\begin{frame}{Early Use of Regions}
Capability Calculus \footcite{crary_typed_1999}
    \vspace{-0.1in}
    \begin{itemize}
        \item Arbitrarily-ordered region allocation/deallocation, via capability tracking
        \item Capability: set of regions presently valid to access
    \end{itemize}
\pause
Also see:
    \begin{itemize}
        \item RC \footcite{gay_language_2001}
        \item Reaps \footcite{berger_reconsidering_2002}
        \item Control-C \footcite{kowshik_ensuring_2002} and Type Homogeneity \footcite{dhurjati_memory_2003}
   \end{itemize}
\end{frame}

\note[itemize]{
\item Capability Calculus:
\item  > Add aliasing info to capabilities; track copies to know when truly revoked
\item  > capability: indicates set of regions presently valid to access
\item  > Arbitrarily-ordered region allocation/deallocation, via capability tracking
\item  > capability: set of regions presently valid to access
\item  > Goal: realistic region support for C, but don't give tons of safety guarantees: more to help
\item  > b/c they allow arbitrary region intermingling, they need reference counts
\item RC:
\item > Region-support for C, focus on preventing dangling pointers % aliasing!
\item > Maintain a reference count of external pointers pointing to objects in region % deleteregion fails if count is zero
\item > Type system + annotations ($\texttt{sameregion}, \texttt{traditional}, \texttt{parentptr}$) to help compiler remove annotation operations % no updating associated region % enforced by runtime checks
\item > Annotations ($\texttt{sameregion}, \texttt{traditional}, \texttt{parentptr}$) to prevent updating associated region % enforced by runtime checks
\item > Restrict C by disallowing arbitrary integer-to-pointer casts
\item > Break cyclic data structures before deleting regions
\item > No memory safety guarantee
\item > 27\%-11\% overhead
\item Reaps:
\item > Combine regions and heaps into ``reaps''; regions until $\texttt{reapFree}$ deallocates individual object
\item > Place freed object on a heap; subsequent allocations use heap until exhausted
\item > Use Lea allocator generally, unless you need fast regions (i.e. reaps) % what their study found
\item Control-C
\item > Single region active at any given time
\item > Automatic Pool Allocation and Type Homogeneity for safe dangling pointers; objects represented by the same points-to graph node are allocated in a common pool; can individually deallocate
\item > Pointer value containing region address must be provably dead before \texttt{rfree}
}

\begin{frame}{Regions in Cyclone}
    Cyclone\footcite{grossman_region-based_2002}, again!
    \begin{itemize}[<+->]
        \item Three region types:
            \begin{itemize}
                \item single \alert{heap} region \note{(lives forever)}
                \item \alert{stack} regions \note{(local-declaration blocks)}
                \item \alert{dynamic} regions \note{(lexically-scoped lifetimes)}
            \end{itemize}
        \item Lifetime subtyping: region A <: region B $\Leftrightarrow$ region A \alert{outlives} region B
        \item Sane defaults by inferring region annotations on pointer types
    \end{itemize}
    \vspace{0.1in}
\end{frame}
\note[itemize]{
\item heap region: lives forever, can be garbage collected
\item lexically-scoped: i.e. region x { int i<x> = 5; ... }
\item Region polymorphism
\item Annotate pointers with regions into which they point; must be live for dereference
\item Can specify partial orders on region lifetimes between pointers
\item intraprocedural analysis => need function prototypes
}

\begin{frame}[fragile]{Regions in Cyclone}
\footnotesize
  \begin{columns}[T]
    \begin{column}{0.52\textwidth}
      \begin{itemize}[<+->]
        \item Pointer can escape scope of their regions \note{via closures/existential types; RC dealt with this via reference counting}
        \item Type system tracks \alert{capability} (live regions set)
        \item Check needed \alert{liveness} on pointer dereference % no closures, but existential types can hide pointers to escape scope of their regions
        \item Function \alert{effect}: set of regions it might access
        \item Calculate function's effect from prototype alone
        %\item Type of elements not known, so use $\epsilon$ to abstract set of regions it might need to access
        %\item Fix with static operator $\texttt{regions\_of}$ %instead of effect variables
      \end{itemize}
    \end{column}

\begin{column}{0.46\textwidth}

    \pause

\begin{lstlisting}[language=C,numbers=left,mathescape,basicstyle={\footnotesize\ttfamily}]
struct Set<$\alpha$, $\rho$, $\epsilon$>
{
  list_t<$\alpha$, $\rho$> elts;
  int (*cmp)($\alpha$,$\alpha$; $\epsilon$);
}
\end{lstlisting}

\pause

\begin{lstlisting}[language=C,numbers=left,mathescape,basicstyle={\footnotesize\ttfamily}]
struct Set<$\alpha$, $\rho$>
{
  list_t<$\alpha$, $\rho$> elts;
  int (*cmp)($\alpha$,$\alpha$; regions_of($\alpha$));
}
\end{lstlisting}

\end{column}
\end{columns}

\end{frame}

\begin{frame}{Some Cyclone Abstract Syntax}
\scriptsize{
\begin{alignat*}{2}
\text{kinds}\quad &\kappa &&::= T \altm R
\\
\text{type and region vars}\quad &\alpha,\rho &&
\\
\text{region sets}\quad &\epsilon &&::= \emptyset \altm \alpha \altm \epsilon_1 \cup \epsilon_2
\\
\text{region constraints} \quad &\gamma &&::= \emptyset \altm \colorIt{red}{n1}{\gamma, \epsilon <: \rho}
\\
\text{constructors} \quad &\tau &&::= \alpha \altm \text{int} \altm \tau_1 \xrightarrow[]{\epsilon} \tau_2 \altm \tau @ \rho \altm \text{handle}(\rho) \altm \forall\alpha:\kappa \triangleright \gamma.\tau \altm ...
\\
\text{expressions} \quad &e &&::= x_\rho \altm v \altm e\ \langle\tau\rangle \altm \ast e \altm \textbf{new}(e_1)e_2 \altm e_1(e_2) \altm \&e \altm ...
\\
\text{functions} \qquad &f &&::= \colorIt{blue}{n2}{\rho:(\tau_1\ x_{\rho}) \xrightarrow[]{\epsilon} \tau_2 = \{s\}} \altm \Lambda\alpha:\kappa \triangleright \gamma.f
\\
\text{statements} \qquad &s &&::= e \altm s_1;s_2 \altm \textbf{if } (e)\ s_1 \textbf{ else } s_2 \altm \rho:\{\tau\ x_\rho = e;\ s\} \altm \textbf{region}\langle\rho\rangle\ x_{\rho}\ s \altm ...
\\
...
\end{alignat*}
    }
\end{frame}

\begin{frame}{Example Cyclone Judgments}
  \begin{columns}[T]
    \begin{column}{0.32\textwidth}
     $\Delta;\Gamma;\gamma;\epsilon;\tau \vdash_{stmt} s$
    \end{column}
    \begin{column}{0.32\textwidth}
     $\Delta;\Gamma;\gamma;\epsilon \vdash e\ :\ \tau$
    \end{column}
    \begin{column}{0.32\textwidth}
     $\Gamma \vdash \epsilon \Rightarrow \rho$
    \end{column}
  \end{columns}

\pause

$\Delta: \text{in-scope type/region vars}$
\pause
\\
$\Gamma: \text{mapping of in-scope vars to types}$
\pause
\\
$\gamma: \text{constraints relating lifetimes}$
\pause
\\
$\epsilon: \text{capability, i.e. live regions}$
% Delta holds type and region variables in scope
% Gamma holds value variables in scope and types
% \gamma records partial order constraints relating region lifetimes
% \epsilon records capability (which regions in \Delta are live)
% tau is type e must have in any statement of form return e
\pause

% expressions access memory can be proven to be live from epsilon and gamma
\begin{center}
\begin{tabular}{cc}
\infer[(\textsc{var})]
{\Delta;\Gamma;\gamma;\epsilon \vdash x_\rho : \Gamma(x_\rho)}
{\gamma \vdash \epsilon \Rightarrow \rho}
    &
\infer[(\textsc{deref})]
{\Delta;\Gamma;\gamma;\epsilon \vdash \ast e:\tau}
{\Delta;\Gamma;\gamma;\epsilon \vdash e:\tau \ast \rho \qquad \gamma \vdash \epsilon \Rightarrow \rho}
\end{tabular}

\pause
\vspace{0.1in}

% every alpha and rho in epsilon_1 can be proven live from epsilon and gamma
\begin{tabular}{c}
\infer[(\textsc{call})]
{\colorIt{blue}{n1}{\Delta;\Gamma;\gamma;\epsilon \vdash e_1(e_2):\tau}}
{\colorIt{red}{n2}{\Delta;\Gamma;\gamma;\epsilon \vdash e_1 : \tau_2 \xrightarrow[]{\epsilon_1} \tau}
    \qquad
    \colorIt{green}{n3}{\Delta;\Gamma;\gamma;\epsilon \vdash e_2 : \tau_2}
    \qquad
    \colorIt{orange}{n4}{\gamma \vdash \epsilon \Rightarrow \epsilon_1}}
\end{tabular}

\pause
\vspace{0.1in}

% type instantiation
\begin{tabular}{c}
\infer[(\textsc{type-inst})]
{\Delta;\Gamma;\gamma;\epsilon \vdash e\langle\tau_1\rangle: \tau_2[\tau_1 / \alpha]}
{\Delta;\Gamma;\gamma;\epsilon \vdash e : \forall\alpha:\kappa\triangleright\gamma_1 . \tau_2
    \qquad
 \Delta \vdash \tau : \kappa
    \qquad
 \gamma \vdash \gamma_1[\tau_1 / \alpha]}
\end{tabular}

% TODO review this
% novelty: ensuring gamma establishes constraints \gamma_1 used when type-checking e
% judgement \gamma \vdash \gamma' means for every \epsilon <: \rho in \gamma', we can show \gamma \vdash \rightarrow \epsilon
% \tau_2[\tau_2/alpha] is capture-avoiding substitution of \tau_1 for \alpha in \tau
% \gamma_1[\tau_1/alpha] is substitution of regions_of(\tau_1) for \alpha in \gamma_1

\end{center}
\end{frame}


\begin{frame}{Cyclone Results}
Soundness Theorem \footcite{grossman_formal_2001}:
    \begin{itemize}
        \item The program cannot get stuck from type errors or dangling-pointer dereferences. 
        %\item If starting statement is well-formed with respect to initial heap, the program cannot get stuck from type errors or dangling-pointer dereferences. 
        \item The terminating program deallocates all regions it allocates \note{non-vacuous -- it's saying b/c it did its job right}
    \end{itemize}
    % can compile with GC
\pause
Benchmarks
    \begin{itemize}
        \item 86 lines of region annotation-related changes across 18,000 lines (6\%) \note{change 10\% of code)}
        \item Eliminate heap allocation entirely for web-server
        \item Near-zero overhead (from garbage collection and bounds checking)
    \end{itemize}
\end{frame}


\subsection{Linear Types and Ownership}

% introduce
\begin{frame}{Linear Types}
    Linear types \footcite{girard_linear_1987}\footcite{wadler_linear_1990} ensure that every variable is used exactly \emph{once}.
    \begin{itemize}[<+->]
        %\item Cannot be arbitraily copied or discarded \note{i.e. allow exchange, not weakening or contraction}
        \item The \alert{world} is a non-duplicatable resource % don't want to world states at same time
        \item Track proper modification of world % (modelled as an array)
        \item Efficiency: safe to destructively update an array
        \item $\textbf{Memory management: can immediately collect used values}$ % only one reference, can show memory management does not leak
        %\item Many read-only references (non-linear), only one reference (linear) for writing %not at same time; write must be only reference
        %\item Paper shows how linear types can determine when an array's update operation can re-use storage in a safe way
    \end{itemize}
\end{frame}
% mention girard's linear logic

%\begin{frame}{Linear Types}
%\begin{alignat*}{2}
%    \text{Exprs}\quad &e &&::= c \altm x \altm \texttt{fn}_{\pi}\ x\ \Rightarrow e_0 \altm e_1\ e_2 \altm \texttt{new}_{\pi}\ x\ :=\ e_1\ \texttt{in}\ e_2 \altm!x \altm x := e_0
%\\
%    \text{Types}\quad &\tau &&::= \texttt{int} \altm \texttt{bool} \altm \tau_1 \rightarrow \tau_2 \altm \tau\ \texttt{ref}
%\end{alignat*}
%\end{frame}

\begin{frame}{Applied Linear Types}
\footnotesize{
Clay \footcite{hawblitzel_low-level_2004}
  \begin{itemize}
    \item Type-theoretic basis for giving type-safe code more control over memory
    \item Singleton types to type check loads, coercion functions to modify values' type safely
    %\item Memory is single linear array words, linearity prevents aliasing
    %\item Type sequences to safely introduce non-linearity
    %\item Used to implement several GCs
  \end{itemize}

\pause

$\textsc{PACLANG}$ \footcite{ennals_linear_2004} %: packet processing
  \begin{itemize}
    \item Program network processors for handling packets % with linear type system
    \item Unique ownership property: each packet in heap is referenced by exactly one thread
    \item Allow mutable aliasing within the same thread
    \item Operations for a functions to 1) take ownership or 2) create local aliases
  \end{itemize}
}
\end{frame}

\begin{frame}{COGENT}
    $\textsc{cogent}$ \footcite{amani_cogent:_2016}\footcite{oconnor_cogent:_2016}
    \begin{itemize}[<+->]
        \item Pure, polymorphic language with linear types for writing low-level systems code \note{high-level; filesystems}
        \item Efficient machine code with in-place updates
        %\item Imperative C semantics and functional value semantics \note{the linear types; for equational reasoning}
        \item Linear types:
            \begin{itemize}
                \item Ensure safe handling of heap-allocated objects
                \item Equational functional semantics via mutable state/imperative effects
            \end{itemize}
        \item Reason with \alert{interactive theorem prover}
        \item Missing functionality can be implemented in C, manually verified
        \item No trusted compiler, runtime, or garbage collector needed
    \end{itemize}
\end{frame}
% total functions (no recursion)
% no recursion, borrowing, mutable state

\begin{frame}{COGENT}
  \begin{columns}[T]
    \begin{column}{0.52\textwidth}
    $\vcenter{\hbox{\includegraphics[height=18em]{fig/FIGURE_oconnor-2016_cogent-chain.PNG}}}$
    \end{column}
\pause
    \begin{column}{0.46\textwidth}
\footnotesize
\vspace{0.2in}
Person-Months of Work
    \begin{itemize}
        \item Proof Framework: 33.5
        \item Compiler: 10
        \item Proofs: 18 \note{design, formalisation, and proof of Cogent and properties}
    \end{itemize}
Lines (kLOC)
    \begin{itemize}
        \item Isabelle theorems: 17
        \item Compiler: 9.5 \note{of Haskell $\textsc{cogent}$ compiler code}
        \item ext2 Filesystem: 6.5 (Isabelle/HOL: 76.7) \note{cogent}
    \end{itemize}
    \end{column}
  \end{columns}
\end{frame}

\note[itemize]{
\item Their verification chain is extensive.
\item Given a cogent program, the cogent compiler first generates C code.
\item A C parser then generates the semantics of this C code in Isabelle/Simpl (*A*).
\item Using the Simple code, the /AutoCoress/ tool generates a monadic functional program (*B*).
\item Then cogent also produces the following representations:
\item 1. monomorphic A-normal deep embedding of the cogent program (*C*)
\item 2. a polymorphic A-normal deep embeeded of the same (*D*)
\item 3. an A-normal shallow embedding (*E*)
\item 4. a 'neat' shallow embedding (syntactically close to cogent compiler input) (*F*)
\item The chain of verification involves proving that the Simpl semantics (*A*) /refines/ the
\item >  monadic semantics (*B*), which /refines/ the update and value semantics of the monomorphic
\item >  deep embedding (*C*), which refines the value semantics of the polymorphic embedding (*D*),
\item >  which refines the HOL that is equal to the shallow embedding (*E*) and (*F*).
\item Finally, user-supplied abstract datatypes (implemented in C) are verified according to the shallow embedding (*F*).
\item The only unverified link the chain is output of the C-to-Isabelle parser that becomes Simpl code (*A*).
\item  kinds to determine is type can be freely shared or discarded
}

\begin{frame}{Other Linear Types}
  Quasi-linear types \footcite{kobayashi_quasi-linear_1999}
    \vspace{-0.09in}
    \begin{itemize}
        \item Distinguish consumed values from those that may be returned
        \item Use $\kappa$ to control how often a variable of type $\tau^{\kappa}$ is used (many times locally)
    \end{itemize}
\pause
  Vault \footcite{deline_enforcing_2001}
    \vspace{-0.09in}
    \begin{itemize}
        \item Keys associate static capabilities with run-time resources %, held-key set, type guards
        \item Annotate functions with effect clause (pre- and post-conditions on held-key set)
        \item Windows 2000 locking errors, IRP ownership model % corresponds to tracked types
    \end{itemize}
\pause
  Ordered types for memory layout \footcite{petersen_type_2003}
    \vspace{-0.09in}
    \begin{itemize}
        \item Variables must be used in order $\Rightarrow$ memory locations \note{restrict linear types by removing exchange property}
        \item \emph{Orderly lambda calculus} for size-preserving memory operations \note{Coercions to manipulate ordered variables in frontier (combine/split to treat as different types)}
    \end{itemize}
  \vspace{0.2in}
\end{frame}

\begin{frame}{Ownership}
\footnotesize
   Types can represent \alert{ownership} and prevent \emph{aliasing} and \emph{mutation} on the same location.
   \\
   \vspace{0.1in}
   \pause
  LCL \footcite{evans_static_1996}
    \vspace{-0.1in}
    \begin{itemize}
        \item $\texttt{owned}$ annotation to denote reference with obligation to release storage
        \item $\texttt{dependent}$ annotation for sharing; user ensures lifetimes contained properly % lifetime of dependent is within lifetime of owner
    \end{itemize}
  \pause
  Ownership Types \footcite{clarke_ownership_1998}
    \vspace{-0.1in}
    \begin{itemize}
      %\item Each object 1) owns a context and 2) is owned by a surrounding context \note{context: nested store partition; program begins in root context}
      \item Object's definition includes \alert{unique} object context that owns it
      %\item Owning an object and having a reference to it are not the same
      %\item Type + context declaration = \alert{ownership type}
    \end{itemize}
  \pause
  Singularity \footcite{fahndrich_language_2006}
    \vspace{-0.1in}
    \begin{itemize}
        \item Type system tracks resources, passes ownership of arguments to callee \note{possibly permanently}
    \end{itemize}
  \vspace{0.1in}
\end{frame}
\note[itemize]{
  \item Clarke: two variables having ownership types with different declared object contexts cannot refer to same part of store, are not aliases
}

\begin{frame}[fragile]{Rust \footcite{matsakis_rust_2014}\footcite{levy_ownership_2015}\footcite{jung_rustbelt:_2017}}
\footnotesize
    \begin{columns}
    \begin{column}{0.46\textwidth}

    \begin{itemize}[<+->]
        \item Ownership and lifetimes \note{every value has uniquer owner}
        \item Type system enforces that objects have unique \alert{owners} \note{functions transfer ownership, shallow copy, no data dupl of vec}
        \item Objects may be borrowed for no longer than owner \note{unique mutable borrows, or multiple immutable borrows, to solve aliasing}
        \item Object automatically \alert{deallocated} when owner leaves scope
        \item Unsafe sections for mutating raw pointers and aliased state \note{raw: no alias tracking, mutual exclusion, communication between threads}
        \item Regions $\approx$ lifetimes (region capabilities $\approx$ lifetime tokens)
    \end{itemize}

    \vspace{0.1in}
    \end{column}

    \pause

    \begin{column}{0.46\textwidth}
\begin{lstlisting}[
    language=Rust,numbers=left,mathescape,basicstyle={\scriptsize\ttfamily},
    moredelim={**[is][{\btHL<8>}]{@1}{@}},
    moredelim={**[is][{\btHL<9>}]{@2}{@}},
    moredelim={**[is][{\btHL<10>}]{@3}{@}},
    moredelim={**[is][{\btHL<11>}]{@4}{@}},
    moredelim={**[is][{\btHL<12>}]{@5}{@}},
    moredelim={**[is][{\btHL<13>}]{@6}{@}},
    moredelim={**[is][{\btHL<14>}]{@7}{@}},
    moredelim={**[is][{\btHL<15>}]{@8}{@}},
]
@1let (snd, rcv) = channel();@
@2join(move || {@
@3      let mut v = Vec::new();@
@3      v.push(0);@
@4      snd.send(v);@
@8      v.push(1);@
    },
@5    move || {@
@6      let v = rcv.recv().unwrap();@
@7      println!("Received: {:?}", v);@
    });
\end{lstlisting}
\end{column}
\end{columns}

\end{frame}
\note[itemize]{
\item performing certain actions on a resource (reading and writing a memory location) should require a 'right' 
\item or capability that is uniquely owned by one alias at any point during the execution of the program.
\item although right is uniquly owned, it cn be transferred from one alias to another
\item can be shared in a controlled manner
\item note it's affine (can drop but not duplicate)
\item Rust-belt: as long as only unsafe code in well-typed rust is confined to libraries satisfying verification conditions, program is safe
\item Rust-belt: tells you verification condition
}

\begin{frame}{Quick Temporal Recap}
    \begin{itemize}[<+->]
      \item Capabilities and pointer-based metadata \note{can leverage info from spatial techniques to get full temporal safety}
      \item Effects and regions \note{look toward effect types as an alternative to GC}
      \item Linear types \note{deallocate after use}
      \item Ownership and borrowing \note{unique owners help determine who is in charge, prevent data races, responsibility for cleaning}
    \end{itemize}
\end{frame}

\begin{frame}{Conclusion}
    \begin{itemize}[<+->]
        \item Memory errors $\equiv$ \emph{type errors}
        \item \emph{Static} memory management \note{like stack-based approach, important for embedded systems}
        \item Isolate unsafe world
        \item Be reasonable and optimistic \note{dynamic approaches still needed, static has come a long way}
    \end{itemize}
\end{frame}
\note[itemize]{
\item safe vs unsafe worlds: ccured, checked-c, rust, effects, cogent
}

\begin{frame}{Thanks}
    \begin{center}
        \Huge{Thanks!}
    \end{center}
\end{frame}

%\appendix

\end{document}

%%% Local Variables:
%%% mode: latex
%%% TeX-master: t
%%% TeX-engine: xetex
%%% End:
